\documentclass[11pt,a4paper]{article}

\usepackage[margin=1in]{geometry}
\usepackage[T1]{fontenc}
\usepackage[utf8]{inputenc}
\usepackage{lmodern}
\usepackage{microtype}
\usepackage{amsmath,amssymb,amsfonts,amsthm}
\usepackage{mathtools}
\usepackage{graphicx}
\usepackage{booktabs}
\usepackage{array}
\usepackage{enumitem}
\usepackage{xcolor}
\usepackage{float}
\usepackage{caption}
\usepackage[colorlinks=true, allcolors=black, breaklinks=true]{hyperref}
\usepackage{authblk}
\usepackage{tikz}
\usetikzlibrary{arrows.meta,positioning,calc,fit}

\captionsetup{font=footnotesize}

\usepackage[
    backend=biber,
    style=authoryear-comp,
    sorting=nyt,
    maxcitenames=2,
    maxbibnames=99,
    giveninits=false,
    uniquename=false,
    uniquelist=false,
    doi=true,
    url=true,
    isbn=false,
    eprint=true
]{biblatex}

\addbibresource{references.bib}
\DeclareDelimFormat{nameyeardelim}{\addspace}
\setlength{\bibhang}{1.5em}
\DeclareFieldFormat[article,inproceedings,incollection]{title}{\mkbibquote{#1}}
\DeclareFieldFormat[book,report,thesis,misc]{title}{\mkbibemph{#1}}
\DefineBibliographyStrings{english}{
  andothers = {et\addabbrvspace al\adddot}
}

\newtheorem{proposition}{Proposition}
\newtheorem{corollary}{Corollary}
\newtheorem{definition}{Definition}
\newtheorem{remark}{Remark}
\newtheorem{assumption}{Assumption}

\newcommand{\X}{\ensuremath{\mathcal{X}}}
\newcommand{\Z}{\ensuremath{\mathcal{Z}}}
\newcommand{\ThetaSpace}{\ensuremath{\Theta}}
\newcommand{\EtaSpace}{\ensuremath{\mathcal{E}}}
\newcommand{\Mnominal}{\ensuremath{\mathcal{M}_{\mathrm{N}}}}
\newcommand{\Gnuis}{\ensuremath{G_{\mathrm{nuis}}}}
\newcommand{\R}{\ensuremath{\mathbb{R}}}
\newcommand{\quotient}{\ensuremath{\mathcal{X}/G_{\mathrm{nuis}}}}
\newcommand{\taumc}{\ensuremath{\tau_{\mathrm{MC}}}}
\newcommand{\taustar}{\ensuremath{\tau_k^\star}}

\title{
\large
\textbf{Fleet-Wide Condition Monitoring via Shared Nominal Representations and Local Radius Calibration}\\
\medskip
Enabling OEM-Centralized, Fleet-Deployed CBM
}

\author[1]{Eduardo Di Santi}
\affil[1]{Department of Applied Mathematics, University of Colorado Boulder
}
\date{}

\begin{document}
\maketitle

%=====================================================================
% Abstract
%=====================================================================
\begin{abstract}

\textbf{Objectives:} Deploying condition monitoring across large fleets of
heterogeneous assets remains challenging because representation learning,
threshold tuning, and fault-labelled commissioning are typically performed
independently for every monitored unit. This work investigates whether these
tasks can be decoupled by sharing a single nominal representation across the
fleet while allowing each asset to identify only its own local operational
radius during commissioning.

\textbf{Methods:} A canonical nominal representation, implemented as a
Conv1D autoencoder, is trained once offline using simulated healthy signals
and transferred unchanged to every asset. Each deployed unit estimates its
operational radius as an empirical residual quantile from a short stream of
unlabelled nominal observations using a convergence-based commissioning
procedure. The methodology is evaluated under controlled operating
conditions using the Case Western Reserve University (CWRU) bearing dataset
and subsequently validated on four independent run-to-failure bearings from
the NASA IMS dataset. Results are compared with complete-data oracle
calibration, direct simulator-threshold transfer, and autoencoders trained
independently on the early-life data of each NASA bearing.

\textbf{Findings:} Direct transfer of the absolute simulator threshold
produced complete failure after simulation-to-field transfer, confirming
that the transferable object is the nominal representation rather than its
absolute reconstruction-error scale. Across all CWRU operating conditions,
convergence-based commissioning achieved 100\% calibration success, an
average post-calibration healthy false-alarm rate of 3.59\%, ball-fault
detection of 99.69\%, and 100\% detection of inner- and outer-race faults,
while recovering operational radii within 2.08\% of the complete-data
oracle. On NASA IMS, independently commissioned assets converged to
asset-specific operational radii after approximately 50 recordings,
identified persistent degradation in all four bearings, and provided an
average warning of 45.4 hours before the documented end of the experiment,
without requiring per-asset representation learning.

\textbf{Novelty:} The contribution is a fleet deployment methodology rather
than a new anomaly-detection architecture. The deployable model is the
shared nominal representation; deployment of a new asset requires only
estimating its local operational radius during commissioning rather than
training a new representation or engineering an asset-specific threshold.

\textbf{Practical Applications:} A single validated nominal representation
can be developed, versioned, and maintained centrally, for example by an
original equipment manufacturer, and deployed unchanged across an operator's
fleet. Each installed asset requires only a short unsupervised
commissioning phase to estimate its local operational radius, eliminating
per-asset representation training and asset-specific threshold engineering.

The deployable model is the shared nominal representation; deployment of a
new asset requires only estimating its local operational radius rather than
training a new model.

\end{abstract}

\noindent\textbf{Keywords:}
condition-based maintenance (CBM),
rail asset management,
fleet-wide deployment,
asset commissioning,
bearing condition monitoring,
OEM deployment,
shared nominal representation,
simulation-to-field transfer,
edge AI,
run-to-failure validation


%=====================================================================
\section{Introduction}
\label{sec:introduction}
%=====================================================================

Modern transportation systems increasingly rely on condition-monitoring technologies to supervise
large populations of operational assets, including rolling-element bearings, axleboxes, traction
motors, gearboxes, pumps, actuators, vehicles, and spatially indexed infrastructure segments
\cite{jardine2006,zhao2019}. Rolling-element bearings are among the most safety-critical of these
components: railway axlebox bearings alone are monitored across thousands of independently
operating vehicles, where defects can propagate to derailment-level failures if undetected
\cite{entezami2020}. Although machine-learning research has produced increasingly capable
anomaly-detection models \cite{chandola2009,zhao2019}, the deployment
problem remains largely unresolved: a detector that performs well for one
asset under one calibrated operating condition does not automatically scale
to hundreds or thousands of independently commissioned units.

In practice, fleet growth therefore scales not only the number of monitored
assets but also the engineering effort required to train, validate, version,
and maintain asset-specific models.

At fleet rollout, several constraints arise simultaneously. First, representative fault examples
are rarely available. Second, waiting months or years for failures to accumulate is operationally
unacceptable. Third, installation-specific factors such as sensor coupling, operating point,
mounting, and environmental noise alter the numerical scale of anomaly scores between assets
\cite{randall2011}. In railway fleets, for example, axlebox bearings of nominally identical rolling
stock are exposed to different track quality, wheel wear, load distribution, and wayside versus
on-board sensor placement, so that a detector calibrated on one vehicle rarely transfers its
absolute decision threshold to another \cite{entezami2020}. Finally, training, calibrating,
validating, and maintaining a separate representation model for every physical unit creates an
engineering burden that grows directly with fleet size.


The practical question is therefore not simply

\begin{quote}
\textbf{Q1.} \emph{How can an anomalous condition be detected?}
\end{quote}

but rather

\textbf{Q2.}
\emph{How can a single nominal representation be deployed across a
heterogeneous fleet while allowing every physical asset to identify only its
own local operational neighbourhood, without per-asset representation
learning or manual threshold engineering?}

This paper addresses Q2 by separating representation engineering from
deployment calibration. A canonical nominal representation is learned once,
offline, from simulated healthy observations and subsequently deployed
unchanged across the fleet. During commissioning, each physical asset
identifies only its own operational radius from a short stream of unlabeled
nominal observations. The shared representation remains fixed throughout the
asset lifetime.

This separation also defines an OEM-centralized deployment model. The shared
nominal representation is engineered, validated, versioned, and maintained
centrally, for example by the original equipment manufacturer (OEM), whereas
each deployed asset performs only a lightweight local commissioning.
Consequently, fleet expansion no longer requires developing,
validating, and maintaining an independent representation model for every
physical unit. Instead, deployment consists only of distributing the shared
representation and estimating one local operational radius for each asset.

From a broader perspective, deployment can be interpreted as a local inverse
problem. The shared representation is recovered once during offline model
engineering, whereas commissioning consists only of identifying the local
operational neighbourhood of each physical asset. In this sense, deployment
does not learn a new representation; it estimates one physical property of
the monitored asset. This interpretation follows the inverse-problem and
quotient-space formulations introduced in our previous work
\cite{disanti2026quotient,disanti2026inverse}, while the present paper
addresses their engineering realization for scalable fleet deployment.

The transferable component is therefore not an anomaly threshold but the
shared nominal representation itself. Absolute reconstruction-error
thresholds generally do not transfer because their numerical scale depends
on the simulator, preprocessing, operating conditions, and residual
statistics. Instead, the proposed methodology transfers one shared nominal
representation together with a fleet-wide nominal coverage, while each asset
independently estimates the local operational radius required to realise
that coverage. The deployable model is therefore common to the fleet; only
the operational radius is asset-specific.

\begin{figure}[htbp]
\centering
\resizebox{\linewidth}{!}{%
\begin{tikzpicture}[
    font=\small,
    >=Latex,
    flow/.style={
        -{Latex[length=2mm]},
        line width=0.55pt
    },
    panel/.style={
        draw,
        rounded corners=2pt,
        line width=0.6pt
    },
    mainbox/.style={
        draw,
        rounded corners=2pt,
        align=center,
        minimum height=0.85cm,
        inner sep=4pt
    },
    assetbox/.style={
        draw,
        rounded corners=2pt,
        align=center,
        minimum width=1.72cm,
        minimum height=0.95cm,
        inner sep=3pt
    },
    paneltitle/.style={
        font=\bfseries,
        align=center
    },
    footnote/.style={
        font=\footnotesize,
        align=center
    }
]

% Fixed canvas: prevents TikZ from creating an oversized bounding box
\path[use as bounding box] (0,0) rectangle (15.4,5.25);

% ============================================================
% LEFT PANEL
% ============================================================
\draw[panel] (0.00,0.15) rectangle (7.70,5.10);

\node[paneltitle] at (3.80,4.72)
    {Conventional per-asset deployment};

\node[
    mainbox,
    text width=5.8cm
] (oem-left) at (3.80,3.90)
    {OEM / engineering organization};

\node[assetbox] (model-a) at (1.30,2.65)
    {Representation\\$M_A$};

\node[assetbox] (model-b) at (3.85,2.65)
    {Representation\\$M_B$};

\node[assetbox] (model-k) at (6.40,2.65)
    {Representation\\$M_K$};

\draw[flow] (oem-left.south) -- (model-a.north);
\draw[flow] (oem-left.south) -- (model-b.north);
\draw[flow] (oem-left.south) -- (model-k.north);

\node[assetbox] (asset-a) at (1.30,1.48)
    {Asset $A$};

\node[assetbox] (asset-b) at (3.85,1.48)
    {Asset $B$};

\node[assetbox] (asset-k) at (6.40,1.48)
    {Asset $K$};

\draw[flow] (model-a.south) -- (asset-a.north);
\draw[flow] (model-b.south) -- (asset-b.north);
\draw[flow] (model-k.south) -- (asset-k.north);

\node[
    footnote,
    text width=8.4cm
] at (3.80,0.55)
    {$K$ independently engineered, validated, versioned,
     and maintained nominal representations};

% ============================================================
% RIGHT PANEL
% ============================================================
\draw[panel] (7.75,0.15) rectangle (15.55,5.10);

\node[paneltitle] at (11.60,4.72)
    {Proposed OEM-centralized deployment};

\node[
    mainbox,
    text width=5.4cm
] (shared) at (11.60,3.62)
    {One shared nominal representation 
     engineered, validated, versioned, and maintained centrally};

\node[assetbox] (local-a) at (9.25,2.00)
    {Asset $A$\\[1mm]
     estimate $\tau_A^\star$};

\node[assetbox] (local-b) at (11.60,2.00)
    {Asset $B$\\[1mm]
     estimate $\tau_B^\star$};

\node[assetbox] (local-k) at (13.95,2.00)
    {Asset $K$\\[1mm]
     estimate $\tau_K^\star$};

\draw[flow] (shared.south) -- (local-a.north);
\draw[flow] (shared.south) -- (local-b.north);
\draw[flow] (shared.south) -- (local-k.north);

\node[
    footnote,
    text width=8.4cm
] at (11.60,0.72)
    {One centrally governed representation plus\\
     lightweight local calibration at each asset};

\end{tikzpicture}%
}

\caption{
Conventional fleet deployment requires engineering, validating,
versioning and maintaining an independent nominal representation for
each monitored asset. The proposed deployment architecture instead
distributes one shared nominal representation unchanged across the
fleet, while each asset estimates only its local operational radius
during commissioning. The NASA IMS experiments demonstrate that these
locally estimated radii differ substantially between physical assets,
despite using the same representation, the same target nominal coverage,
and the same convergence criterion.
}
\label{fig:oem-deployment}
\end{figure}

This formulation also changes the relevant experimental comparison. The
appropriate reference is not another globally shared detector but an
operationally unrealistic upper bound in which every physical asset has its
complete nominal history available for independent representation learning
and calibration. Evaluating Q2 therefore requires answering a third
question:

\begin{quote}
\textbf{Q3.} \emph{What is the operational cost of replacing independently engineered
asset-specific representations with one centrally maintained shared
representation and lightweight local commissioning?}
\end{quote}

Q3 is evaluated in two complementary deployment scenarios. The first uses
the CWRU \cite{smith2015} benchmark to quantify simulation-to-field transfer under controlled
operating-condition variability while comparing the proposed methodology
with ideal asset-specific representations. The second uses the NASA IMS
run-to-failure dataset to evaluate the complete deployment protocol on
independently monitored physical assets undergoing naturally evolving
degradation.

The asset-specific reference methods are trained and calibrated using the complete healthy dataset
of each operating condition. In contrast, the proposed method uses no field observations for
representation learning and estimates the local operational radius from only the first
100 healthy windows, corresponding to approximately ten seconds of signal.

The experimental results show that the cost of scalable deployment is
small. Under controlled CWRU operating conditions, replacing
asset-specific representations with one shared simulator-trained
representation reduced ball-fault detection only marginally (99.69\%
versus 100\%) while preserving complete inner- and outer-race detection.
The local operational radius converged after approximately fifty
commissioning observations without requiring any field retraining.

The NASA IMS experiments further demonstrate that the same deployment
principle extends to independently monitored physical assets. All four
bearings identified distinct local operational radii using the same shared
representation, the same target nominal coverage, and the same convergence
criterion. After commissioning, every bearing subsequently departed
persistently from its own estimated nominal neighbourhood, providing early
warning between 14.5 and 74 hours before the documented end of the
experiment and supporting the interpretation of the operational radius as a
physical property of the local nominal behaviour rather than a tuned
decision threshold.

The contribution of this paper is therefore not a new anomaly-detection
architecture but a deployment methodology for fleet-scale condition
monitoring. Although a Conv1D autoencoder is used as the experimental representation,
the deployment principle is representation-agnostic. The principal contributions are:

\begin{enumerate}[leftmargin=1.5em]
\item an OEM-centralized fleet-deployment architecture in which a single
shared canonical nominal representation is engineered, validated, and
maintained centrally, while deployed assets perform only lightweight local
operational-radius calibration;

\item a direct simulation-to-field deployment methodology in which a single
shared nominal representation is transferred unchanged across heterogeneous
assets without field-data retraining;

\item an online commissioning procedure that estimates each asset's local
operational radius from a short stream of unlabeled observations;

\item a convergence criterion that determines automatically when sufficient
nominal behaviour has been observed to complete commissioning and freeze the
operational radius;

\item experimental evidence that independently monitored physical assets
identify substantially different local operational radii while sharing the
same representation, target nominal coverage, and convergence criterion;

\item an evaluation against ideal asset-specific detectors trained and calibrated with complete
healthy datasets; and

\item an edge-oriented deployment model that eliminates per-asset
representation learning after installation, allowing fleet-wide deployment
through distribution of one shared representation and local unsupervised
calibration.

\end{enumerate}

The remainder of the paper is organized as follows. Section~\ref{sec:related} reviews condition
monitoring, one-class anomaly detection, adaptive calibration, and fleet deployment.
Section~\ref{sec:method} introduces the deployment methodology.
Section~\ref{sec:experimental_design} presents the experimental protocol.
Section~\ref{sec:results} reports the deployment results.
Section~\ref{sec:math_perspective} summarizes the mathematical interpretation underlying the
deployment guarantees.
Section~\ref{sec:discussion} closes with a discussion of industrial implications, limitations,
and future work.

%=====================================================================
\section{Related Work}
\label{sec:related}
%=====================================================================

Condition-based maintenance (CBM) has become a central component of modern
transportation asset management, where condition monitoring supports the
operation of large fleets of heterogeneous assets
\cite{jardine2006,chandola2009,zhao2019}. Considerable progress has been made
in improving anomaly-detection algorithms; however, comparatively little
attention has been paid to how these detectors are engineered, commissioned,
and maintained at fleet scale once thousands of independent physical assets
must be deployed.

Most machine-learning approaches formulate this problem as supervised fault
classification, requiring representative fault examples during training.
While effective on benchmark datasets, these methods assume that deployment
conditions remain close to those observed during model development.

One-class methods, including one-class SVM \cite{scholkopf1999},
neural-network novelty detection \cite{bishop1994}, and autoencoder-based
anomaly detection \cite{hinton2006,sakurada2014}, reduce the dependence on
fault labels by learning only nominal behaviour. More recent
representation-learning approaches further improve anomaly detection by
constructing latent spaces in which healthy operation occupies a compact
region while degraded behaviour becomes more easily separable.

Most existing work therefore focuses on representation learning itself.
Comparatively little attention has been devoted to the deployment
architecture: how a validated nominal representation can be commissioned,
governed, and reused across an entire fleet while allowing each physical
asset to retain its own operational characteristics.

Adaptive thresholds have been widely used to compensate for changing
operating conditions. However, unrestricted adaptation risks absorbing
progressive degradation into the nominal baseline. This motivates separating
calibration from monitoring. In the proposed methodology, adaptation is
limited to a finite cold-start phase during which each asset estimates its
local operational radius. Once convergence is achieved, the shared
representation remains fixed and subsequent departures become diagnostic
evidence rather than input for further adaptation.

The proposed deployment strategy is also related to invariant
representation learning \cite{bronstein2021} and anomaly detection under
distribution shift \cite{li2023invariant}, both of which seek
representations that are insensitive to nuisance variability. Our
contribution, however, lies at a different level. Rather than proposing a
new representation-learning architecture, we address the deployment problem
by combining a shared nominal representation with local operational-radius
calibration, eliminating per-asset representation learning while preserving
a common nominal coverage across the fleet.

A related and rapidly growing literature addresses cross-condition and
cross-machine transfer through domain adaptation, typically aligning source
and target feature distributions via maximum mean discrepancy, adversarial
domain classifiers, or pseudo-labeled fine-tuning
\cite{wang2026fleet,nguyen2025subdomain,chen2026digitaltwin}. These methods
substantially reduce the need for labeled target-domain faults and have
demonstrated strong cross-condition accuracy on CWRU and related
benchmarks. However, they still require access to target-domain
observations during an explicit adaptation or fine-tuning stage, executed
independently for each new operating condition or physical unit. This
reintroduces, at a smaller scale, the same per-asset engineering burden
that fleet-wide deployment seeks to eliminate: every newly installed asset
still triggers a training procedure, even when it requires no fault
labels.

The distinction is therefore not between supervised and unsupervised
deployment, nor between labeled and unlabeled adaptation, but between
architectures that require engineering work whenever a new asset is
deployed and one in which deployment consists only of estimating a local
operational radius while reusing an already validated shared
representation.

From an engineering perspective, this also changes the transferable object.
Existing transfer-learning approaches transfer model parameters and then
adapt them to each target asset. In the proposed methodology, the
transferable engineering artefact is the validated nominal representation
itself, whereas the only asset-specific quantity estimated during
deployment is the local operational radius.

This viewpoint is also consistent with the inverse-problem and
quotient-space formulations introduced in our previous work
\cite{disanti2026inverse,disanti2026quotient}, in which diagnosis is posed
as the recovery of a canonical nominal representation from observations.
The present paper addresses the engineering realization of that framework
for scalable fleet deployment rather than its mathematical formulation.


%=====================================================================
\section{Method}
\label{sec:method}
%=====================================================================

This section formalizes the proposed deployment methodology. The central idea
is to separate representation engineering from deployment commissioning. A
single nominal representation is engineered once and shared across the fleet,
whereas each deployed asset identifies only its own local operational radius
during commissioning.

%---------------------------------------------------------------------
\subsection{Deployment architecture: shared representation, local operational radius}
%---------------------------------------------------------------------

Let $x_{k,t}\in\mathcal{X}$ denote an observation acquired from operational
asset $k$ at time $t$. In the present experiments,
$x_{k,t}$ corresponds to a fixed-length vibration recording.

A shared nominal representation, trained once offline, maps each observation
to a scalar residual score

\begin{equation}
r_{k,t}=r(x_{k,t}).
\label{eq:score}
\end{equation}

During commissioning, each deployed asset estimates its own local
operational radius, denoted by $\tau_k^\star$, from a short sequence of
unlabeled nominal observations. This radius characterizes the local
neighbourhood of the nominal manifold represented by the shared model.
Once commissioning converges, the operational radius is frozen and remains
fixed throughout subsequent monitoring.

An alarm is therefore generated whenever

\begin{equation}
r_{k,t}>\tau_k^\star.
\label{eq:decision}
\end{equation}

The deployment methodology follows three principles:

\begin{enumerate}[leftmargin=1.5em]
\item one centrally engineered, validated, versioned, and governed nominal representation;
\item one lightweight commissioning procedure performed independently by each deployed asset;
\item no representation retraining, fine-tuning, or parameter adaptation after deployment.
\end{enumerate}

This separation directly matches the OEM-centred deployment model introduced
in Section~\ref{sec:introduction}. The shared representation is engineered,
validated, versioned, and maintained centrally by whichever organization
owns the diagnostic model, for example the component OEM, a third-party
analytics provider, or the fleet operator itself. Each deployed asset
performs only a lightweight local commissioning procedure in which its own
operational radius is identified automatically from nominal observations.

Consequently, fleet growth no longer requires engineering additional
representation models. Deploying a new physical asset consists only of
distributing the already validated shared representation and estimating one
local operational radius for that asset.

The transferable engineering artefact is therefore the nominal
representation itself rather than an absolute decision threshold.
Different assets identify different operational radii while sharing exactly
the same representation, the same target nominal coverage, and the same
commissioning procedure.

%=====================================================================
\section{Experimental Design}
\label{sec:experimental_design}
%=====================================================================

The experimental protocol evaluates the proposed deployment methodology rather
than the underlying representation-learning architecture. Across the five
experimental stages, we measure:

\begin{itemize}[leftmargin=1.5em]
    \item false-alarm rate during cold start;
    \item number of observations required for operational-radius convergence;
    \item variability of $\tau_k^\star$ across assets and operating regimes;
    \item post-convergence fault detection performance;
    \item sensitivity to convergence parameters;
    \item agreement between online and complete-data calibration; and
    \item deployment cost avoided by eliminating per-asset representation learning.
\end{itemize}

\paragraph{Stage 1 --- Controlled multi-asset simulation.}

Four bearing streams are generated with a physics-motivated nominal simulator
using 1200-sample windows, a 30\,Hz shaft frequency, and bounded variation in
amplitude, phase, speed, offset, and additive noise.
The shared representation is trained once on 4,000 simulated healthy windows,
with 1,000 additional nominal windows reserved for validation.
Faults are introduced only during evaluation.
A separate synthetic test set contains 750 windows per operating regime.

Each simulated asset begins in nominal operation and independently estimates
its operational radius.
Asset A remains nominal, whereas Assets B--D transition after convergence to
outer-race, inner-race, and ball-fault regimes.

This stage measures

(i) observations required for convergence;

(ii) false alarms before and after convergence;

(iii) separation between reconstruction residuals and the converged operational
radius after fault onset; and

(iv) whether independently deployed assets converge to stable but
asset-specific operational radii.

\paragraph{Stage 2 --- Direct simulation-to-field deployment on CWRU.}

The nominal representation trained in Stage~1 is transferred unchanged to the
Case Western Reserve University (CWRU) Drive-End dataset.
No CWRU observation is used for representation learning, retraining, or
fine tuning.

Each operating condition is treated as an independent deployment environment.
During commissioning, healthy observations are processed sequentially.
Calibration begins after a minimum of $N_{\min}$ healthy windows and continues
until either

(i) the operational radius satisfies the convergence criterion, or

(ii) a maximum calibration budget $N_{\max}$ is reached.

Once convergence is declared, the operational radius is frozen and all
subsequent observations are evaluated using the unchanged shared
representation.

This stage evaluates whether

(i) the learned representation transfers directly from simulation to field;

(ii) local empirical calibration compensates for the non-transferability of
absolute reconstruction-error thresholds; and

(iii) the calibration duration can be determined automatically from the
observed convergence behaviour rather than prescribed a priori.

The four CWRU operating conditions provide a controlled evaluation of
score-scale heterogeneity and deployment calibration rather than four
physically distinct fleet assets.
Validation on independently monitored run-to-failure bearings is addressed in
Stage~4 using the NASA IMS dataset.

\paragraph{Stage 3 --- Reference comparison and calibration ablation.}

To quantify the cost of scalability, we construct an intentionally favorable
but operationally non-scalable reference condition.
For each CWRU operating condition, three detectors are trained and calibrated
independently using the complete healthy dataset:

(i) a Conv1D autoencoder;

(ii) a one-class SVM operating on PCA-reduced standardized windows; and

(iii) an Elliptic Envelope operating on PCA-reduced standardized windows.

Each detector uses the empirical 98th percentile of its own healthy
anomaly-score distribution as its decision boundary.
Fault observations are never used for model fitting or threshold selection.

Using the unchanged simulator-trained representation, we compare three
deployment strategies:

\begin{enumerate}[leftmargin=1.5em]

\item \textbf{Direct threshold transfer:}
apply the absolute simulator threshold unchanged to each CWRU operating
condition;

\item \textbf{Field oracle:}
estimate the operational radius from the complete healthy dataset available
for each operating condition, without modifying the shared representation;

\item \textbf{Proposed online calibration:}
estimate the operational radius sequentially from healthy observations until
the convergence criterion is satisfied or the maximum calibration budget
$N_{\max}$ is exhausted.

\end{enumerate}

This ablation separates representation transfer from calibration transfer.
The representation is identical in all three deployment strategies; only the
information available for estimating the local operational radius changes.

\paragraph{Stage 4 --- Multi-unit run-to-failure validation on NASA IMS.}

A second nominal representation is constructed for the NASA IMS bearing
family using simulated healthy observations matched to the dataset sampling
and operating characteristics.
The representation is trained once and transferred unchanged to each
independently monitored bearing.

Each bearing estimates its operational radius from an early-life nominal
segment using the same convergence-based calibration procedure adopted in the
CWRU experiments.
Once convergence is declared, the operational radius is frozen and monitoring
continues without further adaptation.

This stage evaluates whether the proposed deployment methodology extends from
heterogeneous operating conditions to physically distinct bearing units
undergoing naturally evolving degradation.

We report

(i) observations required for convergence;

(ii) variability of the converged operational radius across bearings;

(iii) healthy-period false-alarm rate;

(iv) first persistent departure after convergence; and

(v) alarm lead time relative to the documented end of the experiment.

\paragraph{Stage 5 --- Edge deployment cost.}

We report model size, parameter count, CPU inference latency for one
1200-sample window, local calibration memory requirements, and the reduction
in transmitted data obtained when only anomaly scores and selected anomalous
windows are communicated instead of the continuous raw signal.


%=====================================================================
\section{Results}
\label{sec:results}
%=====================================================================

The results are organized around three complementary deployment questions.
First, the CWRU experiments establish a favorable but operationally
non-scalable upper reference in which an independent autoencoder,
preprocessing transformation, and complete-data calibration are available
for each operating condition. Second, they quantify whether one
simulator-trained shared representation, combined only with
convergence-based local commissioning, can recover a comparable diagnostic
operating point without field-model retraining. Third, the NASA IMS
run-to-failure experiment evaluates the complete deployment methodology on
independently monitored physical assets undergoing naturally evolving
degradation.

The CWRU evaluation considers four operating conditions corresponding to
1730, 1750, 1772, and 1797\,RPM. Across these conditions, the experiments
contain 1,414 healthy windows, 1,621 ball-fault windows, 1,621
inner-race-fault windows, and 2,846 outer-race-fault windows. All windows
contain 1,200 samples and are subjected to the same evaluation protocol,
while the representation and calibration information available to each
deployment strategy differ according to the experiment.

The NASA IMS evaluation considers four independently monitored bearings from
Experiment~2, each represented by 984 chronologically ordered recordings.
Unlike the controlled CWRU fault-window comparison, this experiment evaluates
whether locally commissioned operational radii remain diagnostically
meaningful throughout complete run-to-failure trajectories.

%---------------------------------------------------------------------
\subsection{Ideal asset-specific reference}
\label{sec:oracle-results}
%---------------------------------------------------------------------

Table~\ref{tab:asset-specific-by-condition} establishes the practical
upper reference obtained when independent autoencoders are trained and
calibrated for each CWRU operating condition using the complete healthy
dataset available for that condition. These experiments
constitute an oracle reference rather than a deployable fleet architecture:
they assume complete nominal data, an independently fitted preprocessing
transformation, and separate representation learning and calibration for
every operating environment.

Each asset-specific autoencoder uses the empirical 98th percentile of its
own healthy reconstruction-error distribution as its decision boundary.
Fault observations are never used for representation learning or threshold
selection.

\begin{table}[H]
\centering
\caption{Ideal asset-specific autoencoder performance on CWRU. Each
autoencoder is trained and calibrated independently using the complete
healthy dataset of the corresponding operating condition. FAR denotes
healthy false-alarm rate and DR denotes fault detection rate.}
\label{tab:asset-specific-by-condition}

\begin{tabular}{lccccc}
\toprule
RPM &
Threshold &
Healthy FAR (\%) &
Ball DR (\%) &
Inner-race DR (\%) &
Outer-race DR (\%)\\
\midrule

1730 &
0.6485 &
2.23 &
100.00 &
100.00 &
100.00 \\

1750 &
0.7590 &
2.23 &
100.00 &
100.00 &
100.00 \\

1772 &
0.7248 &
2.23 &
100.00 &
100.00 &
100.00 \\

1797 &
0.4591 &
2.46 &
100.00 &
100.00 &
100.00 \\

\midrule

Mean &
-- &
$2.29 \pm 0.12$ &
$100.00 \pm 0.00$ &
$100.00 \pm 0.00$ &
$100.00 \pm 0.00$ \\

\bottomrule
\end{tabular}
\end{table}

The asset-specific autoencoders achieve 100\% detection for all evaluated
fault families under every operating condition while reproducing the
expected approximately 2\% healthy false-alarm rate induced by empirical
p98 calibration. This operating point serves as the reference against which
the deployment cost of the proposed methodology is quantified in the
following sections.

The absolute reconstruction-error threshold differs substantially across
operating conditions, ranging from approximately 0.46 to 0.76. This
variation is important for the deployment problem addressed in this paper:
even when independently trained models produce identical diagnostic
performance, their numerical score scales are not interchangeable.

The asset-specific results therefore define a favorable but operationally
non-scalable reference. Each deployment environment receives its own
preprocessing transformation, representation model, and complete-data
calibration. The subsequent experiments quantify the operational cost of
replacing this ideal deployment with one simulator-trained shared
representation commissioned independently at each asset through local
operational-radius estimation, without any field-model training.


%---------------------------------------------------------------------
\subsection{Repeated-run stability of the asset-specific autoencoder}
\label{sec:ae-stability}
%---------------------------------------------------------------------

To assess the effect of stochastic optimization, the complete
asset-specific training and evaluation procedure was repeated five times
using independent random initializations. The architecture,
preprocessing, datasets, training protocol, and empirical p98 calibration
remained unchanged.

Across all operating conditions, healthy false-alarm rate and the
detection rates for ball-, inner-race-, and outer-race faults remained
unchanged across repetitions. Only the numerical reconstruction threshold
varied, as summarized in Table~\ref{tab:ae-threshold-stability}.

\begin{table}[H]
\centering
\begin{tabular}{lc}
\toprule
RPM & Reconstruction threshold \\
\midrule
1730 & $0.649 \pm 0.021$ \\
1750 & $0.759 \pm 0.227$ \\
1772 & $0.725 \pm 0.148$ \\
1797 & $0.459 \pm 0.026$ \\
\bottomrule
\end{tabular}
\caption{Stability of the reconstruction threshold}
\label{tab:ae-threshold-stability}

\end{table}

Although the detector was retrained five times on exactly the same healthy
dataset, the resulting reconstruction thresholds differed because of
stochastic optimization, while the corresponding diagnostic operating
points remained unchanged.

This observation has an important implication for fleet deployment.
Diagnostic performance is associated with the learned nominal
representation rather than with a unique numerical residual scale.
Different optimization trajectories produce diagnostically equivalent
representations while assigning different absolute reconstruction values to
them. Consequently, transferring an absolute threshold between models is
generally not meaningful, whereas transferring the shared nominal
representation and allowing each deployed asset to identify its own local
operational radius is consistent with the observed behaviour.

%---------------------------------------------------------------------
\subsection{Direct simulation-to-field deployment}
\label{sec:shared-deployment-results}
%---------------------------------------------------------------------

The shared autoencoder is trained once using simulated healthy observations
and transferred unchanged to all four CWRU operating conditions. No CWRU
observation is used to retrain, fine-tune, or otherwise modify either the
representation or its preprocessing transformation.

Three calibration conditions are evaluated:

\begin{enumerate}[leftmargin=1.5em]

    \item \textbf{Direct threshold transfer}, in which the absolute empirical
    p98 residual threshold estimated from simulated nominal data is applied
    unchanged to each field operating condition;

    \item \textbf{Field oracle}, in which the empirical p98 operational radius
    is computed retrospectively from the complete healthy field dataset of
    each operating condition, without retraining the shared representation;

    \item \textbf{Proposed online calibration}, in which the local empirical
    p98 operational radius is estimated sequentially from the commissioning
    residual stream, assumed nominal during initialization. Convergence is
    evaluated only after a minimum calibration size $N_{\min}=40$, and
    calibration continues until the stability criterion is satisfied or the
    maximum budget $N_{\max}=200$ is reached.
\end{enumerate}

Once convergence is detected, the corresponding operational radius is frozen
and used for subsequent monitoring. The effective calibration duration is
therefore determined independently for each operating condition rather than
fixed in advance.

Table~\ref{tab:shared-deployment} reports the complete CWRU deployment
results.

\begin{table}[H]
\centering
\caption{Simulation-to-field deployment of the shared autoencoder on CWRU.
FAR denotes healthy false-alarm rate and DR denotes fault detection rate.
The field oracle uses every available healthy window, whereas the proposed
online procedure stops automatically when the empirical p98 radius satisfies
the convergence criterion. Healthy FAR is calculated over the complete
healthy dataset to permit direct comparison between calibration strategies.
A dash indicates that sequential calibration is not applicable.}
\label{tab:shared-deployment}

\resizebox{\linewidth}{!}{
\begin{tabular}{llrrrrrrr}
\toprule
RPM &
Calibration &
$\tau$ &
Healthy FAR &
Ball DR &
Inner DR &
Outer DR &
Healthy windows &
Calibration windows \\
&
&
&
(\%) &
(\%) &
(\%) &
(\%) &
&
\\
\midrule

1730 & Direct threshold transfer
& 0.1027 & 100.00 & 100.00 & 100.00 & 100.00
& 404 & -- \\

1730 & Field oracle
& 0.3578 & 2.23 & 99.75 & 100.00 & 99.86
& 404 & -- \\

1730 & Proposed online calibration
& 0.3520 & 3.71 & 99.75 & 100.00 & 100.00
& 404 & 49 \\

\addlinespace

1750 & Direct threshold transfer
& 0.1027 & 100.00 & 100.00 & 100.00 & 100.00
& 404 & -- \\

1750 & Field oracle
& 0.3434 & 2.23 & 100.00 & 100.00 & 100.00
& 404 & -- \\

1750 & Proposed online calibration
& 0.3440 & 1.98 & 100.00 & 100.00 & 100.00
& 404 & 62 \\

\addlinespace

1772 & Direct threshold transfer
& 0.1027 & 100.00 & 100.00 & 100.00 & 100.00
& 403 & -- \\

1772 & Field oracle
& 0.3512 & 2.23 & 100.00 & 100.00 & 100.00
& 403 & -- \\

1772 & Proposed online calibration
& 0.3550 & 1.24 & 100.00 & 100.00 & 100.00
& 403 & 58 \\

\addlinespace

1797 & Direct threshold transfer
& 0.1027 & 100.00 & 100.00 & 100.00 & 100.00
& 203 & -- \\

1797 & Field oracle
& 0.3886 & 2.46 & 98.03 & 100.00 & 100.00
& 203 & -- \\

1797 & Proposed online calibration
& 0.3869 & 3.45 & 98.52 & 100.00 & 100.00
& 203 & 49 \\

\bottomrule
\end{tabular}}
\end{table}

Direct transfer of the absolute simulator threshold fails under every
operating condition. Although it produces an apparent 100\% detection rate
for all evaluated fault families, it simultaneously classifies every healthy
field window as anomalous. The absolute residual scale learned in simulation
therefore does not transfer directly to the field.

In contrast, local calibration restores an operationally meaningful decision
boundary without modifying the shared representation. The complete-data field
oracle reproduces the expected approximately 2\% healthy false-alarm rate and
achieves 99.45\% average ball-fault detection, 100\% inner-race detection,
and 99.96\% outer-race detection.

The proposed online calibration reaches a closely comparable operating point
using only the observations available during commissioning. All four
operating conditions satisfy the convergence criterion before the maximum
calibration budget is exhausted, requiring 49, 62, 58, and 49 windows for
1730, 1750, 1772, and 1797\,RPM, respectively. The mean calibration duration
is therefore $54.5 \pm 6.6$ windows.

Across the four conditions, the proposed procedure obtains a mean healthy FAR
of $2.60 \pm 1.18$\%, 99.57\% average ball-fault detection, 100\%
inner-race detection, and 100\% outer-race detection. These results are
close to the complete-data oracle despite requiring neither complete prior
knowledge of the healthy field distribution nor independent representation
learning for each operating condition.

The central result is therefore not that the online procedure outperforms
the complete-data oracle. Rather, it approximates the oracle operating point
through a short, automatically terminated commissioning phase while retaining
one fixed fleet-wide representation.

%---------------------------------------------------------------------
\subsection{Cost of scalable deployment}
\label{sec:deployment-cost}
%---------------------------------------------------------------------

Table~\ref{tab:oracle-vs-proposed} compares the proposed shared deployment
with an ideal but operationally non-scalable asset-specific reference. The
comparison intentionally favors the baseline: the oracle detector is trained
and calibrated using the complete healthy dataset of each operating
condition, whereas the proposed methodology performs no field-model training and estimates the local operational radius through convergence-based commissioning,
 requiring only 49–62 healthy observations before deployment.

\begin{table}[H]
\centering
\caption{Ideal but non-scalable asset-specific deployment versus the proposed fleet-scalable
deployment. Values are mean $\pm$ standard deviation across the four CWRU operating
conditions.}
\label{tab:oracle-vs-proposed}
\resizebox{\linewidth}{!}{
\begin{tabular}{lcccccc}
\toprule
Deployment strategy &
Healthy FAR (\%) &
Ball DR (\%) &
Inner DR (\%) &
Outer DR (\%) &
Field model training &
Healthy field data used \\
\midrule

Asset-specific AE oracle
& $2.29 \pm 0.12$
& $100.00 \pm 0.00$
& $100.00 \pm 0.00$
& $100.00 \pm 0.00$
& One model per condition
& Complete dataset \\

Shared AE with field-oracle radius
& $2.29 \pm 0.12$
& $99.20 \pm 1.60$
& $100.00 \pm 0.00$
& $100.00 \pm 0.00$
& None
& Complete dataset \\

Proposed shared AE with online radius
& $3.09 \pm 2.22$
& $99.69 \pm 0.62$
& $100.00 \pm 0.00$
& $100.00 \pm 0.00$
& None
& 100-window warm-up \\

\bottomrule
\end{tabular}}
\end{table}

Replacing four independently trained autoencoders with one simulator-trained
shared model produces no observed loss for inner- or outer-race detection
and less than one percentage point of average difference for ball-fault
detection. The principal cost appears in healthy false alarms, whose mean
increases from 2.29\% to 3.09\%. Most of this increase is concentrated in the
1797\,RPM condition (5.91\%), whereas the remaining conditions range between
0.50\% and 2.97\%.

These results support the central deployment claim of this paper. A single
simulator-trained canonical representation, combined with short local
operational-radius calibration, achieves near-oracle detection performance
without asset-specific representation learning, substantially reducing
deployment complexity.

\begin{table}[H]
\centering
\caption{Agreement between the converged online operational radius and the
complete-data field-oracle p98 reference. Relative difference is computed as
$100|\tau_{\mathrm{online}}-\tau_{\mathrm{oracle}}|/
\tau_{\mathrm{oracle}}$. Post-calibration FAR is evaluated only on healthy
windows occurring after the declared convergence point.}
\label{tab:online-oracle-agreement}

\begin{tabular}{lccccc}
\toprule
RPM &
$\tau_{\mathrm{oracle}}$ &
$\tau_{\mathrm{online}}$ &
Relative difference (\%) &
Calibration windows &
Post-calibration FAR (\%) \\
\midrule
1730 & 0.3578 & 0.3520 & 1.64 & 49 & 3.94 \\
1750 & 0.3434 & 0.3440 & 0.19 & 62 & 1.75 \\
1772 & 0.3512 & 0.3550 & 1.06 & 58 & 0.87 \\
1797 & 0.3886 & 0.3869 & 0.44 & 49 & 3.90 \\
\midrule
Mean & -- & -- & 0.83 & $54.5 \pm 6.6$ & 2.62 \\
\bottomrule
\end{tabular}
\end{table}

%---------------------------------------------------------------------
\subsection{Operational-radius convergence}
\label{sec:radius-convergence-results}
%---------------------------------------------------------------------

Figure~\ref{fig:tau-convergence} illustrates the sequential estimation of
the local operational radius under the fixed shared representation. During
commissioning, the simulator contributes neither an absolute
reconstruction-error threshold nor a standardized residual scale. The only
fleet-wide calibration quantity is the target nominal coverage,
$1-\alpha=0.98$.

For operating condition $k$, the operational radius available after
observing $t$ commissioning residuals is estimated directly as

\begin{equation}
\tau_k(t)
=
\widehat{Q}_{k,t}(0.98),
\label{eq:tau-online-p98}
\end{equation}

where $\widehat{Q}_{k,t}(0.98)$ denotes the empirical 98th percentile of the
local reconstruction residuals observed up to commissioning instant $t$.
The shared representation remains unchanged throughout this process; only
the asset-specific decision boundary is estimated locally.

Convergence is evaluated after a minimum calibration period
$N_{\min}=40$. Using the criterion defined in
Eq.~\eqref{eq:convergence}, calibration terminates when the relative change
in $\tau_k(t)$ remains below the prescribed tolerance for the required
number of consecutive updates. If convergence is not observed, calibration
continues until the maximum commissioning budget $N_{\max}=200$ is reached.
An asset that does not converge within this interval is not released into
active monitoring and is instead assigned a calibration-failure status.

For retrospective comparison, the complete-data field-oracle radius is

\begin{equation}
\tau_{k}^{\mathrm{oracle}}
=
\widehat{Q}_{k,\mathrm{all}}(0.98),
\label{eq:tau-field-p98}
\end{equation}

where the quantile is computed using every available healthy field window
for operating condition $k$. This quantity is unavailable during
commissioning and is used only to assess how closely the online procedure
approximates the operating point obtained with complete nominal knowledge.

All four CWRU operating conditions satisfy the convergence criterion before
the maximum calibration budget is exhausted. Convergence is detected after
49, 62, 58, and 49 windows for 1730, 1750, 1772, and 1797\,RPM,
respectively, corresponding to a mean commissioning duration of
$54.5 \pm 6.6$ windows. At convergence, the resulting operational radii
differ from their complete-data field-oracle references by only 1.64\%,
0.19\%, 1.06\%, and 0.44\%, respectively.

\begin{figure}[H]
\centering
\includegraphics[width=\linewidth]{cwru_tau_convergence.pdf}
\caption{Sequential estimation of the empirical p98 operational radius
during CWRU commissioning. For each operating condition, the blue segment
shows the online estimate $\tau_k(t)=\widehat{Q}_{k,t}(0.98)$ until the
convergence criterion is satisfied. The red dotted line marks the declared
convergence point, at which the operational radius is frozen and active
monitoring begins. The gray segment is a retrospective continuation obtained
by incorporating additional observations known from the benchmark to be
healthy; it is shown only to characterize the subsequent behaviour of the
estimator and is not used to modify the deployed decision boundary. The
horizontal dashed line indicates the complete-data field-oracle p98
reference.}
\label{fig:tau-convergence}
\end{figure}

The retrospective trajectories show that an empirical upper quantile can
continue to change after an apparently stable interval when additional
high-residual nominal observations arrive. This behaviour is expected
because, during short commissioning periods, the empirical p98 estimate is
determined by only a small number of observations in the upper tail. The
post-convergence trajectories therefore do not represent continued
adaptation of the deployed detector; they show how the estimator would have
evolved had calibration remained active.

The operationally relevant result is the close agreement between the frozen
online radius and the complete-data oracle. Despite using only 49--62
commissioning observations, the online procedure recovers an acceptance
boundary within 2\% of the radius obtained retrospectively from all
available healthy observations. The resulting post-calibration healthy
false-alarm rates range from 0.87\% to 3.94\%, remaining close to the
nominal 2\% operating point despite the substantially reduced calibration
information.

The purpose of the convergence criterion is not to prove that the population
98th percentile has been recovered exactly. It provides an operational
stopping rule indicating that the local radius has remained sufficiently
stable for deployment under the selected parameters
$(N_{\min},N_{\max},L,\varepsilon_C,m)$. These parameters may be adjusted to
reflect the commissioning duration and false-alarm tolerance of a particular
application.

The present experiment assumes that commissioning observations are
representative of nominal operation. Failure to obtain a stable operational
radius within $N_{\max}$ may indicate insufficient commissioning data,
non-stationary operating conditions, or an asset whose initial behaviour is
inconsistent with the assumed nominal state. Distinguishing among these
causes is outside the present evaluation, but the calibration-failure status
provides a conservative mechanism that prevents an unconverged asset from
entering active monitoring.

%---------------------------------------------------------------------
\subsection{Robustness to local initialization}
\label{sec:initialization-robustness}
%---------------------------------------------------------------------

A practical deployment question is whether convergence-based local
calibration depends critically on the particular nominal observations
encountered during commissioning. To evaluate this sensitivity, the shared
simulator-trained representation, preprocessing transformation, target
coverage, and convergence parameters were kept fixed throughout the
experiment. Only the nominal observations available during commissioning,
together with their arrival order, were varied.

For each CWRU operating condition, five independent commissioning sequences
were generated by randomly permuting the available healthy observations.
Calibration was allowed to start after $N_{\min}=40$ observations and
continued until the empirical p98 operational radius satisfied the fixed
convergence criterion or the maximum commissioning budget
$N_{\max}=200$ was exhausted. Once convergence was declared, the estimated
operational radius was frozen and subsequently used for fault detection.
Healthy false-alarm rate was evaluated only on healthy observations excluded
from calibration, while all fault classes were evaluated using the same
shared representation without retraining.

Table~\ref{tab:initialization-robustness} summarizes the results obtained
across the five independent commissioning sequences.

\begin{table}[H]
\centering
\caption{Robustness of convergence-based local calibration to commissioning
sample variability. The shared simulator-trained representation and all
calibration parameters remained fixed; only the healthy commissioning
observations and their arrival order were varied. Values are reported as
mean $\pm$ standard deviation over five independent commissioning
realizations. FAR is evaluated only on healthy observations excluded from
calibration.}
\label{tab:initialization-robustness}

\resizebox{\linewidth}{!}{
\begin{tabular}{lccccccc}
\toprule
RPM &
$\tau_{\mathrm{online}}$ &
Relative oracle error (\%) &
Post-calibration FAR (\%) &
Ball DR (\%) &
Inner DR (\%) &
Outer DR (\%) &
Calibration windows \\
\midrule

1730 &
$0.3486 \pm 0.0023$ &
$2.59 \pm 0.63$ &
$5.38 \pm 1.39$ &
$99.95 \pm 0.11$ &
$100.00 \pm 0.00$ &
$100.00 \pm 0.00$ &
$55.6 \pm 8.2$ \\

1750 &
$0.3451 \pm 0.0111$ &
$2.53 \pm 1.65$ &
$4.23 \pm 5.63$ &
$99.90 \pm 0.14$ &
$100.00 \pm 0.00$ &
$100.00 \pm 0.00$ &
$57.4 \pm 7.9$ \\

1772 &
$0.3511 \pm 0.0086$ &
$1.82 \pm 1.37$ &
$2.54 \pm 2.18$ &
$100.00 \pm 0.00$ &
$100.00 \pm 0.00$ &
$99.97 \pm 0.06$ &
$54.4 \pm 5.9$ \\

1797 &
$0.3915 \pm 0.0065$ &
$1.40 \pm 1.02$ &
$2.21 \pm 2.23$ &
$97.98 \pm 0.71$ &
$100.00 \pm 0.00$ &
$99.97 \pm 0.06$ &
$64.4 \pm 26.0$ \\

\bottomrule
\end{tabular}}
\end{table}

All 20 commissioning runs satisfied the convergence criterion before
reaching the maximum calibration budget, yielding a calibration success
rate of 100\%. Across all operating conditions, the estimated operational
radius remained within $2.08 \pm 0.71$\% of the complete-data oracle, while
the fleet-averaged post-calibration healthy false-alarm rate was
$3.59 \pm 1.59$\%. Fault-detection performance remained essentially
unchanged throughout all commissioning realizations, indicating that the
shared representation is insensitive to realistic variability in the local
initialization process.

These experiments evaluate the robustness of the complete commissioning
procedure rather than the variability of a fixed-size percentile estimate.
The representation, preprocessing pipeline, target coverage, convergence
criterion, and all deployment parameters remained unchanged. Only the local
healthy observations available during commissioning and their arrival order
were modified.

Some variability is expected because the empirical p98 operational radius is
estimated from a relatively small number of upper-tail observations.
Consequently, different commissioning samples naturally produce slightly
different estimates of the local operational radius. These differences
translate into modest variations in the healthy false-alarm rate while
leaving the resulting diagnostic operating point essentially unchanged.

The largest healthy false-alarm rate observed during the repeated
commissioning experiments was 13.8\%. However, this isolated realization did
not alter the deployment outcome: convergence was still achieved, the
estimated operational radius remained within 4\% of the oracle estimate,
and fault-detection performance remained essentially unchanged.

The observed variability therefore reflects finite-sample uncertainty in
estimating an extreme empirical quantile rather than instability of the
proposed deployment methodology. From the quotient-space formulation, this
behaviour is expected: the operational radius estimates the local
neighbourhood of the nominal manifold rather than representing an arbitrary
decision threshold. Consequently, different healthy commissioning samples
should produce slightly different neighbourhood estimates while preserving
the underlying nominal representation.

Together with the oracle-agreement results of
Table~\ref{tab:online-oracle-agreement}, these experiments establish that
convergence-based commissioning reliably recovers the operational
neighbourhood of the shared nominal manifold under realistic initialization
variability. The oracle-agreement experiment evaluates chronological
calibration against the complete-data reference, whereas the present
experiment demonstrates that essentially the same neighbourhood is recovered
under different commissioning realizations without modifying the shared
representation.

The independent NASA IMS validation presented in
Section~\ref{sec:nasa-results} extends this conclusion beyond controlled
operating conditions. There, independently commissioned physical assets
converged to distinct local operational radii while subsequently exhibiting
persistent departures from their respective nominal neighbourhoods during
naturally evolving degradation. Taken together, the CWRU and NASA
experiments support the central deployment hypothesis of this work: the
transferable object is the shared nominal representation, whereas each asset
identifies only its own local operational radius during commissioning.

%---------------------------------------------------------------------
\subsection{Multi-unit run-to-failure deployment on NASA IMS}
\label{sec:nasa-results}
%---------------------------------------------------------------------

The CWRU experiments demonstrate simulation-to-field transfer under controlled
operating conditions. The NASA IMS experiments evaluate the complete deployment
methodology under realistic fleet conditions, where independently operating
physical assets undergo naturally evolving degradation over their remaining
service life.

Each NASA bearing is treated as an independent monitored asset. The
simulator-trained autoencoder and preprocessing pipeline are transferred
unchanged, and no NASA observations are used for representation learning,
retraining, or parameter adaptation. The only asset-specific operation is the
estimation of the local operational radius from only the early-life recordings selected for commissioning and their arrival order were varied.

Table~\ref{tab:nasa-main-results} summarizes the deployment results.
Operational-radius calibration converged successfully for all four bearings
after only 49--57 healthy recordings. Although the estimated operational
radii differ substantially across bearings, every asset uses exactly the same
shared representation, the same target nominal coverage (p98), and the same
convergence criterion without bearing-specific retraining.

Following commissioning, all four bearings subsequently departed persistently from their estimated nominal neighbourhood.
The learned operational radius remains predictive throughout the subsequent degradation trajectory, with a mean
post-departure exceedance rate of 94.6\% and lead times ranging from
14.5~h to 74.0~h before the documented end of the experiment.

\begin{table}[H]
\centering
\caption{NASA IMS deployment results using one shared simulator-trained
representation and convergence-based local operational-radius estimation.}
\label{tab:nasa-main-results}

\resizebox{\linewidth}{!}{
\begin{tabular}{lccccc}
\toprule
Bearing &
Calibration
(windows) &
$\tau_k$ &
Relative
error (\%) &
Post-departure
exceedance (\%) &
Lead
time (h)
\\
\midrule

1 &
49 &
0.010680 &
0.10 &
99.12 &
74.00
\\

2 &
57 &
0.014950 &
0.35 &
86.42 &
38.83
\\

3 &
49 &
0.018930 &
2.87 &
95.88 &
14.50
\\

4 &
49 &
0.006672 &
0.18 &
97.02 &
54.33
\\

\midrule

Mean &
51 &
-- &
0.87 &
94.61 &
45.42
\\

\bottomrule
\end{tabular}}
\end{table}

These results demonstrate that the transferable engineering artefact is the shared nominal representation rather than an asset-specific detector. The same representation identifies substantially different operational radii for different physical assets while requiring no representation retraining and no threshold engineering. Once identified during commissioning, each operational radius remains predictive throughout the subsequent degradation trajectory. The same learned
representation produces substantially different operational radii for
different physical assets while requiring neither representation retraining
nor threshold engineering. Once estimated during commissioning, the local
operational radius continues to separate nominal from degrading behaviour
throughout the remaining lifetime of the asset.

%---------------------------------------------------------------------
\subsubsection{Stability of local operational-radius estimation}
%---------------------------------------------------------------------

A practical deployment methodology must be robust to the particular healthy
observations collected during commissioning. To evaluate this property, the
shared representation was kept completely fixed while only the chronological
commissioning sample used to estimate the local operational radius was varied.

For each NASA bearing, five independent commissioning realizations were generated by sampling and permuting the first 200 early-life recordings.
No retraining was performed; only the observed nominal residuals differed between runs.

Across twenty independent commissioning experiments (four bearings and five
random initialization samples), calibration converged successfully in every
case. Table~\ref{tab:nasa-stability} summarizes the resulting variability.

\begin{table}[H]
\centering
\caption{Stability of local operational-radius estimation under repeated
NASA commissioning experiments using one fixed shared representation.}
\label{tab:nasa-stability}

\resizebox{\linewidth}{!}{
\begin{tabular}{lcccccc}
\toprule
Bearing &
$\tau_{\rm mean}$ &
$\tau_{\rm std}$ &
Relative error (\%) &
Healthy FAR (\%) &
Calibration
(windows) &
Success
(\%)
\\
\midrule

1 &
0.010590 &
0.000074 &
0.73 $\pm$ 0.70 &
5.84 $\pm$ 3.79 &
49.2 $\pm$ 0.4 &
100
\\

2 &
0.014873 &
0.000084 &
0.36 $\pm$ 0.44 &
3.71 $\pm$ 3.33 &
49.0 $\pm$ 0.0 &
100
\\

3 &
0.018151 &
0.000383 &
1.58 $\pm$ 1.88 &
4.93 $\pm$ 4.22 &
52.2 $\pm$ 7.2 &
100
\\

4 &
0.006640 &
0.000033 &
0.66 $\pm$ 0.49 &
5.45 $\pm$ 2.39 &
49.6 $\pm$ 1.3 &
100
\\

\midrule

Mean &
-- &
-- &
0.83 &
4.98 &
50.0 &
100
\\

\bottomrule
\end{tabular}}
\end{table}

These experiments demonstrate that the operational radius is simultaneously
asset-specific and statistically stable. Different physical units identify
different local neighbourhoods of the same shared nominal representation,
whereas repeated commissioning of the same unit produces only minor variation
in the estimated radius. The transferable engineering artefact is therefore
the shared nominal representation itself, while commissioning consists only
of identifying the local operational neighbourhood associated with each
physical asset.

%---------------------------------------------------------------------
\subsection{Edge deployment implications}
\label{sec:edge-results}
%---------------------------------------------------------------------

The shared model contains 847,601 trainable parameters and is trained entirely offline.
Deployment at the monitored asset requires only forward inference, a small residual buffer,
estimation of an empirical quantile, and the current convergence state.
No gradient computation, model optimization, fault labels, or asset-specific representation
learning are required after installation.

On the evaluation platform, inference over a 1,200-sample window requires a median CPU latency
of only 6.26\,ms (95th percentile: 6.62\,ms), substantially below the 100\,ms acquisition
interval corresponding to the evaluated windows. Consequently, the proposed deployment operates
comfortably in real time while leaving ample computational margin for additional edge tasks such
as communication, diagnostics, or signal preprocessing.

The local deployment state is also extremely compact. Besides the shared model itself, each asset
stores only a buffer of the most recent 200 reconstruction residuals during commissioning, requiring less than 1\,kB of asset-specific state. The only
fleet-wide artifact is the shared nominal representation, which is trained once and distributed
to every operational unit.

\begin{table}[H]
\centering
\small
\caption{Edge-deployment footprint of the proposed shared nominal representation.}
\label{tab:edge}
\begin{tabular}{lc|lc}
\toprule
Metric & Value & Metric & Value \\
\midrule
Model parameters & 847,601 & Residual-buffer memory & 800 bytes \\
Raw float32 size & 3.39 MB & Total calibration state $<1$ kB \\
Serialized (.keras) & 10.24 MB & Per-asset model training & None \\
CPU latency (median) & 6.26 ms & Fault labels required & None \\
CPU latency (p95) & 6.62 ms & Local residual buffer & 200 samples \\
\bottomrule
\end{tabular}
\end{table}


\section{Mathematical Scope and Deployment Guarantees}
\label{sec:math_perspective}

This section states, without full proof, the minimal functional-analytic scaffolding justifying why one shared canonical representation and independently estimated local operational radii are mathematically compatible; complete derivations are given in our companion technical reports \cite{disanti2026quotient,disanti2026inverse}.

\begin{quote}
\textbf{Scope.} Jointly recovering a nuisance-invariant representation \emph{and} calibrating per-asset decision boundaries across a heterogeneous fleet is combinatorially hard. This paper decouples that problem into (i) a \emph{global} inverse problem, solved once offline, recovering a canonical representation whose nominal orbit is relatively compact (Proposition~\ref{prop:compactness}); and (ii) a \emph{local} inverse problem, solved independently by each asset as an $O(1)$-per-observation empirical-quantile estimate.
\end{quote}

\subsection{Observation model}
\label{sec:obs_model}

Let $X \subset \mathbb{R}^n$ denote observation windows, $\theta \in \Theta$ the latent diagnostic state, and $\eta \in E$ nuisance variables (sensor coupling, mounting, operating point, noise), with $x = F(\theta,\eta)$. The group $G_{\text{nuis}}$ acts on $E$ while preserving $\theta$; observations are diagnostically equivalent, $x \sim x'$, iff $x' = F(\theta, g\cdot\eta)$ for some $g \in G_{\text{nuis}}$.

\subsection{Compactness and factorization}
\label{sec:compactness}

\begin{proposition}[Relative compactness of the nominal orbit]
\label{prop:compactness}
If $E$ is compact and $\eta \mapsto F(\theta_N,\eta)$ is uniformly bounded and equicontinuous on $E$, then $O_N = \{F(\theta_N,\eta):\eta\in E\}$ is relatively compact in $(C(X), \|\cdot\|_\infty)$.
\end{proposition}
\begin{proof}[Proof sketch]
By Arzel\`a--Ascoli \cite{rudin1976}, boundedness and equicontinuity on compact $E$ give relative compactness of $O_N$; the equicontinuity modulus for the Stage~1 simulator is verified in \cite{disanti2026inverse}.
\end{proof}

Because the nominal orbit is compact, $\tau_k^\star$ only accounts for asset-specific scale around a fixed structure rather than searching an unbounded space, explaining the rapid convergence observed empirically ($51.0\pm12.6$ windows on CWRU, Section~\ref{sec:radius-convergence-results}).

\begin{proposition}[Factorization through the nuisance quotient]
\label{prop:factorization_prop}
Let $\pi: X \to X/G_{\text{nuis}}$ be the quotient projection. If $d:X\to\mathbb{R}$ is $G_{\text{nuis}}$-invariant, then $d$ is constant on fibers and factors uniquely as $d = \bar{d}\circ\pi$.
\end{proposition}
\begin{proof}[Proof sketch]
Invariance makes $d$ constant on fibers, so $\bar{d}([x]):=d(x)$ is well defined by the universal property of the quotient map \cite{disanti2026inverse}.
\end{proof}

Together, these results justify treating diagnosis as membership testing around a fixed canonical representation: $\bar{d}$ (equivalently $O_N$) is recovered once, globally, while each asset $k$ solves only a well-posed local calibration, estimating $\tau_k^\star$ for coverage $1-\alpha$ (Eq.~\eqref{eq:adaptive}) from its own residuals. Empirically, this is reflected in the same fixed representation transferring across four CWRU conditions with only $\tau_k(t)$ re-identified ($99.69\%$ average ball-fault, $100\%$ inner/outer-race detection), and extending to physically distinct NASA IMS bearings with $83.25$~h average lead time.

\section{Discussion}
\label{sec:discussion}

Calibration should not be interpreted as parameter tuning. Rather, it identifies an unknown physical property of the monitored asset: its nominal operational radius.

The principal result of this work is not that a Conv1D autoencoder can separate CWRU fault
classes, a capability already demonstrated extensively in the literature. The contribution lies
instead in the deployment methodology: a single canonical nominal representation can be shared
across heterogeneous assets, while each physical unit independently estimates its own operational
radius without retraining the representation.

From a procurement perspective, this separation gives fleet operators a concrete basis for
specifying deployment requirements: rather than accepting per-asset models developed ad hoc by
each vendor, an operator can require a single validated, versioned nominal representation
together with documented local calibration behaviour, independent of which organization,
component OEM, third-party analytics provider, or in-house engineering team, produces and
maintains it.

The repeated-run and initialization-robustness experiments jointly support this claim from two
complementary directions. Independent random initializations converge to slightly different
reconstruction-error scales, yet healthy false-alarm rate and fault detection remain essentially
unchanged after empirical calibration. Conversely, keeping the learned representation fixed while
varying only the commissioning sample produces only modest variation in the estimated operational
radius. The transferable object is therefore not an absolute reconstruction-error threshold but a
shared canonical representation together with a target nominal coverage. The calibration-transfer
experiments demonstrate that the representation transfers successfully from simulation to field,
whereas the absolute residual scale does not.

Identifying the shared representation as the transferable object also changes the engineering
organization of fleet deployment. Rather than developing and maintaining an independent
representation model for each physical asset, these activities can be centralized, for example by
the original equipment manufacturer (OEM), while each deployed asset performs only a lightweight
local calibration during commissioning. Computational scalability is correspondingly decoupled
from fleet size: growth increases the number of monitored assets, not the number of engineered
models, since adding a unit requires only distributing the shared representation and estimating
its local operational radius.

Beyond commissioning, the NASA IMS experiments indicate that the estimated operational radius
remains predictive throughout the subsequent degradation trajectory. Across all evaluated
bearings, calibration converged after approximately fifty healthy observations, while persistent
degradation subsequently evolved outside the learned nominal neighbourhood for the vast majority
of the remaining asset lifetime. Although the four bearings estimated substantially different
operational radii, all of them used the same shared representation, the same target nominal
coverage (p98), and the same convergence criterion, without bearing-specific retraining.

This behaviour is consistent with the broader theoretical framework developed in our previous
work, which formulates industrial diagnosis as an inverse problem over quotient spaces
\cite{disanti2026inverse,disanti2026quotient,disanti2025blueprints}. Within that formulation,
the operational radius is not merely a tuned threshold but an estimate of the local neighbourhood
of the nominal manifold. Consequently, once calibration has converged, trajectories that remain
nominal are expected to remain inside this neighbourhood, whereas persistent degradation should
naturally evolve outside it. The NASA IMS experiments provide empirical evidence consistent with
this theoretical prediction.

This perspective shifts commissioning from threshold tuning to manifold identification. Rather
than adjusting detector parameters for each individual asset, deployment consists of estimating
the local operational neighbourhood associated with a shared nominal representation. The learned
representation becomes the transferable engineering artefact, while commissioning reduces to
identifying the operational radius of each physical asset.

Operational-radius convergence also clarifies the role of adaptation: its purpose is to identify
the asset during commissioning, not to track its aging indefinitely. Once $\tau_k^\star$ has
been established, continued movement of the anomaly score constitutes diagnostic evidence rather
than input for further adaptation. Initialization can therefore proceed until a predefined
stability criterion is satisfied, after which the operational radius is frozen or only tightly
constrained. Together with the edge footprint reported in
Section~\ref{sec:edge-results}, this deployment model enables day-one monitoring without fault
labels, isolates asset-specific adaptation to local calibration, and requires only lightweight
edge computation.

Conventional CBM deployments typically require asset-specific engineering of both models and
alarm thresholds. The present results suggest that this engineering effort can instead be
replaced by a standardized commissioning procedure in which every asset estimates its own
operational radius while reusing a common nominal representation.

\begin{remark}
The variability observed for the 1750\,RPM commissioning samples is not
specific to the shared-representation architecture. It reflects the
finite-sample uncertainty inherent in estimating an upper empirical
quantile before the complete nominal distribution is available. An
asset-specific detector calibrated from the same limited commissioning
sample would be subject to the same statistical limitation; the
complete-data oracle avoids it only by assuming retrospective access to all
healthy observations.
\end{remark}

The convergence rule adopted in the present experiments is operational
rather than universal. Calibration is declared after the relative
variation of the empirical p98 radius remains below
$\varepsilon_C=0.01$ over a look-back interval $L=10$ for
$m=10$ consecutive updates. These parameters define a configurable trade-off
between commissioning duration and robustness against transient plateaus in
the upper-tail estimate. Requiring more consecutive stable updates, a
longer look-back interval, or a tighter tolerance would generally delay
deployment while reducing sensitivity to short-term fluctuations;
conversely, less restrictive settings would shorten commissioning at the
cost of increased threshold variability. The present study does not claim
that $(L,\varepsilon_C,m)=(10,0.01,10)$ is optimal. Rather, it demonstrates
that this single fixed configuration achieved successful convergence for
all evaluated CWRU and NASA IMS commissioning experiments while preserving
near-oracle diagnostic behaviour. Systematic selection of these parameters
according to application-specific false-alarm and commissioning-time
requirements remains an important direction for future work.

\subsection{Limitations}
\label{sec:limitations}

Several limitations qualify the scope of the present results and should guide their
interpretation.

\textbf{Nominal initialization assumption.} The procedure assumes the cold-start observations are
representative of healthy operation. If an asset were already degraded at installation,
$\tau_k^\star$ could absorb part of that condition into the nominal baseline; a consistency check
against the canonical nominal manifold is a natural extension and is left as future work.

\textbf{Short-warm-up sensitivity.} Because the operational radius is an empirical 98th-percentile
estimated from only 100 windows, it can be sensitive to a small number of extreme values, as
observed in the 1797\,RPM condition (5.91\% FAR, roughly twice the other conditions). Longer
warm-ups or robust quantile estimators would likely reduce this variability at the cost of a longer
cold start.

\textbf{Stationarity of nominal operation.} The convergence criterion assumes an approximately
stationary healthy residual distribution during cold start. The retrospective trajectories in
Figure~\ref{fig:tau-convergence} show modest post-calibration drift; more variable startup
conditions were not evaluated and may benefit from a longer or adaptively triggered calibration
window.

\textbf{Validation domain.} The validation uses bearing datasets (CWRU, NASA IMS) rather than
instrumented transportation assets. While the deployment principle is architecture- and
modality-agnostic by construction (Section~\ref{sec:method}), its extension to other components
remains to be confirmed empirically.



\subsection{Future Work}

A natural extension relaxes the assumption that assets begin in nominal condition. Because the
methodology separates the shared representation from the local radius, future work will
investigate validating the initialization phase itself against the canonical nominal manifold, to
distinguish genuinely healthy startup from incipient degradation before it is absorbed into the
baseline. A second direction concerns progressive degradation: whether the joint evolution of
reconstruction residuals and operational-radius dynamics carries predictive information,
connecting this methodology with physics-informed health indicators and remaining-useful-life
estimation.

Finally, the mathematical interpretation in \cite{disanti2026quotient} suggests formalizing the
operational radius as an asset-specific neighborhood around the transferable canonical
representation, characterizing its invariance properties and the conditions under which local
calibration preserves the underlying quotient-space structure.


%=====================================================================
\section{Conclusion}
%=====================================================================

We presented a day-one deployment methodology for nominal condition monitoring across
heterogeneous asset fleets, validated on public rolling-element bearing datasets. A single
nominal representation is trained once on simulated healthy observations and transferred
unchanged to field deployment. Each operational unit then performs only a short unlabeled
warm-up to estimate its own operational radius, without modifying or retraining the shared
representation.

The experimental results show that a simulator-trained nominal representation transfers
successfully to field deployment, whereas an absolute decision threshold does not.
A shared target nominal coverage, realized locally through the empirical residual quantile estimated during initialization, 
provides an effective bridge between the shared representation and field deployment. This deployment strategy preserves detection performance
close to that of asset-specific models trained under ideal asset-specific calibration conditions
while eliminating per-asset representation learning.

The transferable object is therefore not an absolute anomaly threshold, nor an asset-specific
representation, but a shared canonical nominal representation together with a lightweight local
calibration mechanism. The central outcome of this work is consequently not a new anomaly
detector, but a scalable deployment principle: learn the canonical nominal representation once,
distribute it across the fleet, and let each physical asset identify only its own operational
radius during installation.

This separation between representation learning and deployment calibration provides a practical
foundation for large-scale edge condition monitoring, where an OEM can centrally produce and
maintain the shared representation while each deployed asset performs its own lightweight local
calibration automatically, without operator intervention or manual threshold tuning.

Conventional CBM deployments typically require asset-specific engineering of models and alarm thresholds. The present results suggest that this engineering effort can instead be replaced by a standardized commissioning procedure in which every asset estimates its own operational radius while reusing a common nominal representation.
The representation is shared. The operational radius is individual.

%\vspace{-0.3cm}
%\setlength{\bibhang}{1em}
%\setlength{\bibitemsep}{0.4em plus 0.02em minus 0.02em}
%\renewcommand{\bibfont}{\footnotesize}
%\enlargethispage{3\baselineskip}
\printbibliography

\end{document}
